\section{Introduction}
\label{sec:intro}


Software engineering courses require student teams to deliver a working product
together with requirements, tests, reports, and evidence of contribution. In
practice, these artifacts are often managed in separated tools. Requirements may be
written in documents, tasks may be tracked in a Kanban board, screenshots may be
stored in chat messages, test cases may be updated late, and code activity may stay
only in GitHub. This makes progress visible at the task level, but not always
verifiable at the requirement level. A leader or mentor may still be unable to
answer whether a requirement has been implemented, tested, supported by accepted
evidence, and included in a report.

Requirement traceability addresses this problem by connecting requirements with
downstream artifacts such as tasks, tests, source code, and change records
\cite{nlpTraceability,crossLevelTraceability}. However, maintaining traceability is
difficult in academic projects because teams work quickly, use lightweight tools,
and often prepare documentation near the end of a semester. At the same time, AI
tools can help draft tasks, tests, and reports, but AI-generated text is not reliable
unless it is grounded in structured project data. A fluent weekly report is not
enough if it cannot be traced back to requirements, evidence, and code activity.

DevTrack AI is proposed as an evidence-based traceability and reporting system for
student software teams. It is not designed as another task board. Instead, it
organizes project work around the artifact flow:

\begin{center}
\texttt{Requirement -> Use Case -> Task -> Test Case -> Bug/Defect -> Evidence -> RTM -> Report}
\end{center}

The goal is to make project progress explainable, not only visible. A task should
not be considered complete simply because its status is moved to \texttt{Done}; it
should be linked to a requirement, checked by tests or review evidence, and visible
in the Requirement Traceability Matrix (RTM). The RTM then becomes a live project
health view rather than a static table written only for final submission.

This paper contributes a lightweight system method for academic project tracking.
First, DevTrack AI uses rule-based RTM status calculation so that requirement health
is transparent. Second, it combines an Evidence Vault with a leader-controlled
review gate to reduce unsupported task completion. Third, it uses GitHub integration
and Code Insight to connect code evidence with project artifacts. Fourth, it uses AI
only as a support layer for suggestions and report drafts, while official decisions
remain under human review.
